% =====================================================================
%  2bat-ud6-11.tex  —  SESSIÓ 11: Llegir i qüestionar un informe
%  (sense preàmbul: s'inclou des de main.tex)
% =====================================================================
\horatitol{Sessió 11 — Llegir i qüestionar un informe}%
{Activitat 4 (fase 1): fer una lectura crítica estructurada d'un informe social real amb dades inferencials.}

\subsection{Idea clau}

Comencem l'activitat que tanca l'etapa: la lectura crítica d'un \textbf{informe social
real} amb dades. En parelles, aplicareu \textbf{tot} el que s'ha treballat per valorar, amb
criteri, la solidesa de les conclusions —no per desconfiar de tot, sinó per saber fins on
són fiables.

\begin{idea}
\textbf{Guió per llegir críticament un informe.} Preguntes que cal fer-se sempre:
\begin{enumerate}[leftmargin=1.6em, itemsep=2pt]
  \item \textbf{Què conclou} l'informe? (la tesi principal)
  \item \textbf{Amb quines dades i quina mostra?} (mida, com s'ha triat)
  \item \textbf{És la mostra representativa} o pot estar esbiaixada?
  \item \textbf{Com expressa la incertesa?} (marges d'error, intervals, o els amaga?)
  \item \textbf{Les diferències que destaca són significatives} o podrien ser casuals?
  \item \textbf{Quines limitacions reconeix} (o amaga) l'estudi?
\end{enumerate}
\textbf{L'objectiu:} valorar la solidesa amb criteri, ni creure-s'ho tot ni desconfiar de
tot per sistema.
\end{idea}

\noindent\textbf{Informe (resum).} Una entitat publica: «\emph{El nostre programa d'inserció
funciona: el $68\%$ dels participants troba feina en sis mesos, molt per sobre de la mitjana
del sector}». Fitxa: $90$ participants d'una sola edició; sense grup de comparació; marge
d'error no indicat; enquesta feta als qui van completar el programa.

\subsection{Exemples}

\begin{exemple}[aplicar el guió a l'informe]
Analitzem l'informe amb el guió:
\begin{itemize}[leftmargin=1.4em, itemsep=2pt]
  \item \textbf{Conclou} que el programa funciona (un $68\%$ troba feina).
  \item \textbf{Mostra:} $90$ persones, una sola edició, i \textbf{només qui va completar}
        el programa (biaix de selecció: qui va abandonar no compta).
  \item \textbf{Incertesa:} no dona marge d'error; amb $90$ persones, el marge és notable.
  \item \textbf{Comparació:} diu «per sobre de la mitjana del sector» però no dona la xifra
        ni un grup de comparació: no es pot valorar si la diferència és significativa.
\end{itemize}
\textbf{Conclusió crítica:} el resultat és \emph{prometedor} però \textbf{no concloent}
tal com es presenta.
\end{exemple}

\begin{exemple}[el que caldria per fer-lo sòlid]
Perquè la conclusió fos sòlida, l'informe hauria d'incloure: una mostra més gran i de
diverses edicions; el resultat \textbf{de qui va abandonar} també (no només dels qui van
completar); el \textbf{marge d'error}; i una \textbf{comparació justa} (un grup similar sense
el programa, o la xifra concreta del sector amb el seu interval). Sense això, «funciona» és
una afirmació poc fonamentada.
\end{exemple}

\subsection{Exercicis resolts}

\begin{exresolt}[detectar el biaix principal]
Quin és el biaix principal de la mostra de l'informe i com afecta la conclusió?
\solucio
El biaix principal és de \textbf{selecció}: s'enquesta \textbf{només qui va completar} el
programa, i no qui el va abandonar. Això \textbf{sobreestima} l'èxit, perquè justament els
que no van encaixar (i potser no van trobar feina) queden fora del càlcul. La conclusió
«funciona» queda inflada.
\resposta{Biaix de selecció (només qui completa); sobreestima l'èxit.}
\end{exresolt}

\begin{exresolt}[qüestionar la comparació]
L'informe diu que el $68\%$ és «molt per sobre de la mitjana del sector». Per què no es pot
donar per bona aquesta comparació tal com està?
\solucio
Perquè no dona la \textbf{xifra} de la mitjana del sector ni el seu \textbf{interval}, ni
diu si els grups són comparables. Sense saber quant és aquesta mitjana i amb quin marge, no
es pot valorar si la diferència és \textbf{significativa} o si podria ser casual. La
comparació, tal com està, no es pot avaluar.
\resposta{Falta la xifra i l'interval del sector; la comparació no es pot avaluar.}
\end{exresolt}

\subsection{Exercicis per fer}

\noindent\textit{Analitzeu l'informe del resum aplicant el guió de lectura crítica.}

\begin{exercicis}
\begin{exlist}
  \item Quina és la conclusió principal de l'informe? Amb quina mostra la sosté?
  \item És representativa la mostra? Identifica'n el biaix principal i digues en quina
        direcció afecta el resultat.
  \item L'informe expressa la incertesa (marge d'error)? Per què és un problema que no ho
        faci?
  \item La comparació amb «la mitjana del sector» es pot valorar? Què hi falta?
  \item \textbf{(Limitacions)} Enumera dues limitacions que l'informe no reconeix
        explícitament.
  \item \textbf{(Conclusió crítica)} Escriu, en tres o quatre línies, una valoració
        global de la solidesa de l'informe: què s'hi pot creure, amb quines cauteles, i què
        caldria per fer-lo concloent.
\end{exlist}
\end{exercicis}

% ---------------------------------------------------------------------
\fullsolucions{Sessió 11}
\begin{solucions}
\begin{sollist}
  \item Conclou que el programa d'inserció funciona ($68\%$ troba feina en sis mesos). Ho
        sosté amb $90$ participants d'una sola edició, enquestats només si van completar el
        programa.
  \item No és representativa: biaix de selecció (només qui va completar). Sobreestima
        l'èxit, perquè exclou qui va abandonar (que potser no va trobar feina).
  \item No dona marge d'error. És un problema perquè, amb $90$ persones, la incertesa és
        notable i el $68\%$ podria variar bastant; presentar-lo com a xifra tancada amaga
        aquesta incertesa.
  \item No es pot valorar: falta la xifra concreta de la mitjana del sector, el seu
        interval i si els grups són comparables. Sense això no se sap si la diferència és
        significativa.
  \item Resposta oberta. Per exemple: (1) una sola edició i mostra petita (poc
        generalitzable); (2) cap grup de comparació que permeti aïllar l'efecte del
        programa; (3) exclusió dels qui van abandonar.
  \item Resposta oberta. Idea: el resultat és prometedor però no concloent; s'hi pot creure
        com a indici favorable, amb la cautela que la mostra és petita, esbiaixada i sense
        marge ni comparació. Per fer-lo sòlid caldria mostra més gran i de diverses
        edicions, incloure qui va abandonar, donar el marge d'error i una comparació justa.
\end{sollist}
\end{solucions}
